Om aan een WiFi netwerk te kunnen connecten moeten we eerst weten of een WiFi netwerk beschikbaar is met voldoende sterkte:
\begin{lstlisting}[language=bash]
$ nmcli device wifi list
\end{lstlisting}

Als het netwerk waaraan we willen connecten beschikbaar is met voldoende signaal sterkte dan kunnen we eraan connecten:
\begin{lstlisting}[language=bash]
$ nmcli device wifi connect <SSID_OR_BSSID> password <PASSWORD>
\end{lstlisting}
Dit maakt een profiel bestand aan in /texttt{/etc/NetworkManager/system-connections/} met de naam van het SSID of BSSID.

Als je wilt connecten aan een verborgen (hidden) network gebruik dan de optie \texttt{hidden}:
\begin{lstlisting}[language=bash]
$ nmcli device wifi connect <SSID_OR_BSSID> password <PASSWORD> hidden yes
\end{lstlisting}

Als je meer dan 1 WiFi interface hebt en met een specifieke WiFi interface wilt verbinden, geef dan de interface naam op:
\begin{lstlisting}[language=bash]
$ nmcli device wifi connect <SSID_OR_BSSID> password <PASSWORD> ifname <WLANIF>
\end{lstlisting}

Met het show commando kunnen we een overzicht krijgen van onze verbindingen:
\begin{lstlisting}[language=bash]
$ nmcli connection show
\end{lstlisting}
De laatste kolom vertelt ons aan welke interface een SSID of BSSID verbonden is.

Om de verbinding af te sluiten gebruiken we:
\begin{lstlisting}[language=bash]
$ nmcli device disconnect ifname <WLANIF>
\end{lstlisting}

Om een netwerk te activeren die we eerder hebben geconfigureert gebruiken we het bestaande profiel (naam of UUID):
\begin{lstlisting}[language=bash]
$ nmcli connection up <NAME_OR_UUID>
\end{lstlisting}

Willen we een profiel weggooien, dus een netwerk volledig vergeten, dan is het \texttt{delete} commando je vriend:
\begin{lstlisting}[language=bash]
$ nmcli connection delete <NAME_OR_UUID>
\end{lstlisting}

En tot slotte kunnen we WiFi ook volledige uit zetten:
\begin{lstlisting}[language=bash]
$ nmcli radio wifi off
\end{lstlisting}

